\subsection{Diagonalization, Commutants, and Simultaneous Forms}

\begin{definition} \label{def:diagonalizable-operator} \hlt{Diagonalizable Operator}\\
Let $\mathcal{V}$ be a finite-dimensional vector space over $\mathbb{F}$. An operator $T\in\LL(\mathcal{V})$ is \hlt{diagonalizable} iff $(\exists\text{ a basis }\beta\text{ of }\mathcal{V})\,[T]_{\beta\leftarrow\beta}$ is diagonal.
\end{definition}

\begin{proposition} \label{prop:eigenspaces-for-distinct-eigenvalues-form-direct-sum} \hlt{Eigenspaces for Distinct Eigenvalues Form a Direct Sum}\\
Let $\mathcal{V}$ be finite-dimensional, let $T\in\LL(\mathcal{V})$, and let $\lambda_1,\ldots,\lambda_m$ be distinct eigenvalues of $T$, where $m\geq1$. The following hold:
\begin{enumerate}[label=(\roman*)]
\setlength{\itemsep}{0pt}
\item The sum $E_{\lambda_1}(T)+\cdots+E_{\lambda_m}(T)$ is direct.
\item $\sum\limits_{j=1}^{m}\dim E_{\lambda_j}(T)\leq\dim(\mathcal{V})$.
\end{enumerate}
\begin{proof}
\begin{enumerate}[label=(\roman*)]
\setlength{\itemsep}{0pt}
\item Let $\mathbf{u}_j\in E_{\lambda_j}(T)$ for $j\in\{1,\ldots,m\}$, and suppose $\sum\limits_{j=1}^{m}\mathbf{u}_j=\mathbf{0}$. Set $J\defeq\{j\in\{1,\ldots,m\}:\mathbf{u}_j\neq\mathbf{0}\}$.\\
If $J\neq\emptyset$, then $(\mathbf{u}_j)_{j\in J}$ is linearly independent by Theorem \ref{thm:distinct-eigenvalues-independent-eigenvectors}, contradicting its zero relation with all coefficients equal to $1$.\\
Thus $J=\emptyset$, and the sum is direct by Proposition \ref{prop:internal-direct-sum-family-criterion}(ii),(i).
\item The identity $\dim(E_{\lambda_1}(T)\oplus\cdots\oplus E_{\lambda_m}(T))=\sum\limits_{j=1}^{m}\dim E_{\lambda_j}(T)$ holds by Theorem \ref{thm:direct-sum-iff-dimensions-add}(i).\\
Its left side is at most $\dim(\mathcal{V})$ by Proposition \ref{prop:dimension-bounds-and-basis-criteria}(i).
\end{enumerate}
\hlt{Structural Commentary}: Eigenvector independence converts a mixed eigenspace relation into the direct-sum criterion.\\
\hlt{Key Insight}: Distinct-eigenvalue independence forbids cross-eigenspace cancellation.
\end{proof}
\end{proposition}

\begin{theorem} \label{thm:conditions-equivalent-to-diagonalizability} \hlt{Conditions Equivalent to Diagonalizability}\\
Let $\mathcal{V}$ be finite-dimensional, $T\in\LL(\mathcal{V})$, $n\defeq\dim(\mathcal{V})$. Let $(\lambda_1,\ldots,\lambda_m)$ enumerate $\operatorname{spec}(T)$, where $m\geq0$. The following are equivalent:
\begin{enumerate}[label=(\roman*)]
\setlength{\itemsep}{0pt}
\item $T$ is diagonalizable;
\item $\mathcal{V}$ has a basis consisting of eigenvectors of $T$;
\item $(\exists\,(\mathcal{U}_j)_{j=1}^{n}\in\operatorname{Inv}(T)^n)$ such that $(\forall j\in\{1,\ldots,n\})\,\dim(\mathcal{U}_j)=1$ and $\mathcal{V}=\bigoplus\limits_{j=1}^{n}\mathcal{U}_j$;
\item $\mathcal{V}=\bigoplus\limits_{j=1}^{m}E_{\lambda_j}(T)$;
\item $\dim(\mathcal{V})=\sum\limits_{j=1}^{m}\dim E_{\lambda_j}(T)$.
\end{enumerate}
\begin{proof}
\begin{enumerate}[label=(\roman*)]
\setlength{\itemsep}{0pt}
\item[$(i)$] $\Rightarrow (ii)$ Choose a basis $\beta=(\mathbf{v}_1,\ldots,\mathbf{v}_n)$ for which $[T]_{\beta\leftarrow\beta}$ is diagonal by Definition \ref{def:diagonalizable-operator}.\\
For $j\in\{1,\ldots,n\}$, the $j$th matrix column has the form $a_j\mathbf{e}_j$, so $T\mathbf{v}_j=a_j\mathbf{v}_j$ by Definition \ref{def:matrix-representation-of-linear-map}. Thus $\beta$ is an eigenbasis.
\item[$(ii)$] $\Rightarrow (iii)$ Let $(\mathbf{v}_1,\ldots,\mathbf{v}_n)$ be eigenbasis, and $\mathcal{U}_j\defeq\operatorname{span}(\mathbf{v}_j)$. Each $\mathcal{U}_j$ is one-dimensional, invariant under $T$ by Proposition \ref{prop:one-dimensional-invariant-subspaces-and-eigenvectors}.\\
The equality $\mathcal{V}=\sum\limits_{j=1}^{n}\mathcal{U}_j$ holds by spanning of the basis, while every zero relation across the summands is a basis relation.\\
Hence the sum is direct by Proposition \ref{prop:internal-direct-sum-family-criterion}(ii),(i).
\item[$(iii)$] $\Rightarrow (iv)$ Choose $\mathbf{0}\neq\mathbf{v}_j\in\mathcal{U}_j$ for $j\in\{1,\ldots,n\}$. Each $\mathbf{v}_j$ is an eigenvector by Proposition \ref{prop:one-dimensional-invariant-subspaces-and-eigenvectors}, so $\mathcal{U}_j$ lies in an eigenspace of $T$.\\
Hence the listed eigenspaces span $\mathcal{V}$. If $m=0$ and $n\geq1$, then $\mathbf{v}_1$ would be an eigenvector, contradicting $\operatorname{spec}(T)=\emptyset$.\\
Thus $m=0$ implies $n=0$, and the asserted empty direct sum is $\mathcal{V}=\{\mathbf{0}\}$. If $m\geq1$, the sum is direct by Proposition \ref{prop:eigenspaces-for-distinct-eigenvalues-form-direct-sum}(i).
\item[$(iv)$] $\Rightarrow (v)$ If $m=0$, condition (iv) says $\mathcal{V}=\{\mathbf{0}\}$, so both sides of (v) are zero. If $m\geq1$, the dimensions add by Theorem \ref{thm:direct-sum-iff-dimensions-add}(i).
\item[$(v)$] $\Rightarrow (i)$ If $m=0$, condition (v) implies $n=0$, and the empty basis diagonalizes $T$. Suppose $m\geq1$.\\
For $j\in\{1,\ldots,m\}$, choose a basis $\beta_j$ of $E_{\lambda_j}(T)$ and let $\beta$ be their concatenation. \\
The eigenspace sum is direct by Proposition \ref{prop:eigenspaces-for-distinct-eigenvalues-form-direct-sum}(i), so $\beta$ is a basis of that sum by Theorem \ref{thm:direct-sum-iff-dimensions-add}(ii).\\
Its length is $n$ by (v), and it is independent, hence it is a basis of $\mathcal{V}$ by Proposition \ref{prop:full-dimension-basis-criteria}(ii).\\
Each member of $\beta$ is an eigenvector, so $[T]_{\beta\leftarrow\beta}$ is diagonal by Definition \ref{def:matrix-representation-of-linear-map}. Therefore $T$ is diagonalizable by Definition \ref{def:diagonalizable-operator}.
\end{enumerate}
\hlt{Structural Commentary}: The implication cycle links diagonal form, eigenbases, and eigenspace dimension.\\
\hlt{Key Insight}: Diagonalizability is exactly exhaustion by distinct eigenspaces.
\end{proof}
\end{theorem}

\begin{corollary} \label{cor:enough-eigenvalues-implies-diagonalizable} \hlt{Enough Eigenvalues Imply Diagonalizability}\\
Let $\mathcal{V}$ be finite-dimensional and let $T\in\LL(\mathcal{V})$. If $\lvert\operatorname{spec}(T)\rvert=\dim(\mathcal{V})$, then $T$ is diagonalizable.
\begin{proof}
Set $n\defeq\dim(\mathcal{V})$. If $n=0$, the empty basis diagonalizes $T$. Suppose $n\geq1$, enumerate $\operatorname{spec}(T)$ as $\lambda_1,\ldots,\lambda_n$, and choose $\mathbf{0}\neq\mathbf{v}_j\in E_{\lambda_j}(T)$ for $j\in\{1,\ldots,n\}$. The list $(\mathbf{v}_1,\ldots,\mathbf{v}_n)$ is linearly independent by Theorem \ref{thm:distinct-eigenvalues-independent-eigenvectors}.\\
Its length is $n$, so it is a basis by Proposition \ref{prop:full-dimension-basis-criteria}(ii). Hence $T$ is diagonalizable by Theorem \ref{thm:conditions-equivalent-to-diagonalizability}(ii),(i).\\
\hlt{Structural Commentary}: Distinct eigenvalues supply an independent eigenvector list whose cardinality reaches the ambient dimension.\\
\hlt{Key Insight}: A full-length independent eigenvector list is already a diagonalizing basis.
\end{proof}
\end{corollary}

\begin{perspective} \label{per:diagonalizability-viewpoint} \hlt{Diagonalization as Exhaustion by Eigenspaces}\\
Separation holds by Proposition \ref{prop:eigenspaces-for-distinct-eigenvalues-form-direct-sum}(i), while spanning is the remaining condition by Theorem \ref{thm:conditions-equivalent-to-diagonalizability}.\\
Diagonalization therefore records exhaustion of $\mathcal{V}$ by its distinct eigenspaces.
\end{perspective}

\begin{theorem} \label{thm:algebraic-spectral-resolution-diagonalizable} \hlt{Algebraic Spectral Resolution}\\
Let $\mathcal{V}$ be nonzero and finite-dimensional, $T\in\LL(\mathcal{V})$, $\Lambda\defeq\{\lambda_1,\ldots,\lambda_m\}\subseteq\mathbb{F}$, where $m\geq1$, $\lambda_i$ are distinct. The following hold:
\begin{enumerate}[label=(\roman*)]
\setlength{\itemsep}{0pt}
\item $T$ is diagonalizable with $\operatorname{spec}(T)=\Lambda$ iff $(\exists\,\rho_1,\ldots,\rho_m\in\LL(\mathcal{V})\setminus\{0\})$ such that $(\forall i)\,\rho_i^2=\rho_i$, $(\forall i,j)\,(i\neq j\Rightarrow\rho_i\rho_j=0)$,\\
$\sum\limits_{i=1}^{m}\rho_i=I_{\mathcal{V}}$, and $T=\sum\limits_{i=1}^{m}\lambda_i\rho_i$, where all indices range over $\{1,\ldots,m\}$.
\item Every family $(\rho_1,\ldots,\rho_m)$ satisfying the right side of part (i) obeys $(\forall i\in\{1,\ldots,m\})\,\operatorname{range}(\rho_i)=E_{\lambda_i}(T)$.
\item Every such family obeys $(\forall i\in\{1,\ldots,m\})\,\ker(\rho_i)=\bigoplus\limits_{j\in\{1,\ldots,m\}\setminus\{i\}}E_{\lambda_j}(T)$.
\item Every such family obeys $(\forall i\in\{1,\ldots,m\})\,\rho_i=\prod\limits_{j\in\{1,\ldots,m\}\setminus\{i\}}\frac{T-\lambda_jI_{\mathcal{V}}}{\lambda_i-\lambda_j}$, where the empty product is $I_{\mathcal{V}}$.
\end{enumerate}
\begin{proof}
\begin{enumerate}[label=(\roman*)]
\setlength{\itemsep}{0pt}
\item[$(i)$] ($\Rightarrow$) The decomposition $\mathcal{V}=\bigoplus\limits_{i=1}^{m}E_{\lambda_i}(T)$ holds by Theorem \ref{thm:conditions-equivalent-to-diagonalizability}(i),(iv).\\
For $i\in\{1,\ldots,m\}$, $\mathcal{W}_i\defeq\bigoplus\limits_{j\in\{1,\ldots,m\}\setminus\{i\}}E_{\lambda_j}(T)$, $\rho_i\defeq\rho_{E_{\lambda_i}(T),\mathcal{W}_i}$. Each $\rho_i$ is algebraic projection by Proposition \ref{prop:algebraic-projection-range-kernel-idempotent}(ii).\\
If $\mathbf{x}=\sum\limits_{k=1}^{m}\mathbf{x}_k$ with $\mathbf{x}_k\in E_{\lambda_k}(T)$, then $\rho_i\mathbf{x}=\mathbf{x}_i$ by the projection formula.\\
Hence $\rho_i\neq0$, $(\forall i\neq j)\,\rho_i\rho_j=0$, $\sum\limits_{i=1}^{m}\rho_i=I_{\mathcal{V}}$ by coordinate action. Thus $T=\sum\limits_{i=1}^{m}\lambda_i\rho_i$ by $T\mathbf{x}=\sum\limits_{i=1}^{m}\lambda_i\mathbf{x}_i=\sum\limits_{i=1}^{m}\lambda_i\rho_i\mathbf{x}$ $\forall \mathbf{x}\in\mathcal{V}$.
\item[$(i)$] ($\Leftarrow$) Fix a family satisfying the right side, and set $\mathcal{V}_i\defeq\operatorname{range}(\rho_i)$ for $i\in\{1,\ldots,m\}$. The equality $\mathcal{V}=\sum\limits_{i=1}^{m}\mathcal{V}_i$ holds by $\sum\limits_{i=1}^{m}\rho_i=I_{\mathcal{V}}$.\\
If $\mathbf{x}_i\in\mathcal{V}_i$ and $\sum\limits_{i=1}^{m}\mathbf{x}_i=\mathbf{0}$, then $\mathbf{x}_k=\rho_k\sum\limits_{i=1}^{m}\mathbf{x}_i=\mathbf{0}$ by idempotence and pairwise annihilation.\\
Hence $\mathcal{V}=\bigoplus\limits_{i=1}^{m}\mathcal{V}_i$ by Proposition \ref{prop:internal-direct-sum-family-criterion}(ii),(i). For $\mathbf{x}\in\mathcal{V}_i$, $T\mathbf{x}=\lambda_i\mathbf{x}$ by idempotence and pairwise annihilation, so $\mathcal{V}_i\subseteq E_{\lambda_i}(T)$. The inclusion $\lambda_i\in\operatorname{spec}(T)$ holds by $\mathcal{V}_i\neq\{\mathbf{0}\}$ and $\mathcal{V}_i\subseteq E_{\lambda_i}(T)$, where $\mathcal{V}_i\neq\{\mathbf{0}\}$ by $\rho_i\neq0$.\\
If $T\mathbf{x}=\mu\mathbf{x}$ with $\mathbf{x}\neq\mathbf{0}$, write $\mathbf{x}=\sum\limits_{i=1}^{m}\mathbf{x}_i$ with $\mathbf{x}_i\in\mathcal{V}_i$. Then $(\forall i)\,(\lambda_i-\mu)\mathbf{x}_i=\mathbf{0}$ by directness, and some $\mathbf{x}_i\neq\mathbf{0}$.\\
Hence $\mu\in\Lambda$, so $\operatorname{spec}(T)=\Lambda$ by the two inclusions. The eigenspaces span $\mathcal{V}$ by the inclusions $\mathcal{V}_i\subseteq E_{\lambda_i}(T)$ and the direct-sum decomposition, so $T$ is diagonalizable by Theorem \ref{thm:conditions-equivalent-to-diagonalizability}(iv),(i).
\item[$(ii)$] Fix family satisfying right side of part (i), set $\mathcal{V}_i\defeq\operatorname{range}(\rho_i)$.\\
Identities $\mathcal{V}=\bigoplus\limits_{i=1}^{m}\mathcal{V}_i$m$\mathcal{V}_i\subseteq E_{\lambda_i}(T)$ hold by reverse-direction calculation in part (i).\\
If $\mathbf{x}\in E_{\lambda_i}(T)$, write $\mathbf{x}=\sum\limits_{j=1}^{m}\mathbf{x}_j$ with $\mathbf{x}_j\in\mathcal{V}_j$. Then $(\forall j\neq i)\,\mathbf{x}_j=\mathbf{0}$ by $\sum\limits_{j=1}^{m}(\lambda_j-\lambda_i)\mathbf{x}_j=\mathbf{0}$, directness, and distinctness.\\
Hence $\mathbf{x}\in\mathcal{V}_i$, and $\mathcal{V}_i=E_{\lambda_i}(T)$ by the two inclusions.
\item[$(iii)$] For $\mathbf{x}\in\mathcal{V}$, write $\mathbf{x}=\sum\limits_{j=1}^{m}\mathbf{x}_j$ with $\mathbf{x}_j\in\mathcal{V}_j$ by the reverse-direction calculation in part (i).\\
Then $\rho_i\mathbf{x}=\mathbf{x}_i$ by idempotence and pairwise annihilation. Hence $\mathbf{x}\in\ker(\rho_i)$ iff $\mathbf{x}_i=\mathbf{0}$, and $\ker(\rho_i)=\bigoplus\limits_{j\in\{1,\ldots,m\}\setminus\{i\}}\mathcal{V}_j$ by this equivalence. The stated formula follows by part (ii).
\item[$(iv)$] For $i\in\{1,\ldots,m\}$, set $\ell_i(z)\defeq\prod\limits_{j\in\{1,\ldots,m\}\setminus\{i\}}\frac{z-\lambda_j}{\lambda_i-\lambda_j}$.\\
If $\mathbf{x}\in E_{\lambda_k}(T)$, then $\ell_i(T)\mathbf{x}=\ell_i(\lambda_k)\mathbf{x}=\delta_{ik}\mathbf{x}=\rho_i\mathbf{x}$ by part (ii), idempotence, and pairwise annihilation.\\
The eigenspaces span $\mathcal{V}$ by the reverse direction of part (i) and part (ii), so $\ell_i(T)=\rho_i$ by equality on a spanning family.
\end{enumerate}
\hlt{Structural Commentary}: Eigenspace decomposition proves the equivalence; interpolation recovers each projection.\\
\hlt{Key Insight}: Lagrange polynomials isolate the individual eigenspaces.
\end{proof}
\end{theorem}

\begin{lemma} \label{lem:diagonalizable-restriction-to-invariant-subspace} \hlt{Diagonalizable Restrictions to Invariant Subspaces}\\
Let $\mathcal{V}$ be nonzero finite-dimensional vector space, $T\in\LL(\mathcal{V})$ be diagonalizable, $\mathcal{U}\in\operatorname{Inv}(T)$.\\
For $\lambda\in\operatorname{spec}(T)$, set $\mathcal{U}_{\lambda}\defeq\mathcal{U}\cap E_{\lambda}(T)$. The following hold:
\begin{enumerate}[label=(\roman*)]
\setlength{\itemsep}{0pt}
\item $\mathcal{U}=\bigoplus\limits_{\lambda\in\operatorname{spec}(T)}\mathcal{U}_{\lambda}$.
\item The restriction $T|_{\mathcal{U}}$ is diagonalizable.
\end{enumerate}
\begin{proof}
For $\lambda\in\operatorname{spec}(T)$, the polynomial spectral projection $\rho_{\lambda}$ exists by Theorem \ref{thm:algebraic-spectral-resolution-diagonalizable}(i),(iv).\\
The inclusion $(\forall k\in\mathbb{Z}_{\geq0})\,T^k(\mathcal{U})\subseteq\mathcal{U}$ holds by induction on $k$ and $T$-invariance. Hence $\rho_{\lambda}(\mathcal{U})\subseteq\mathcal{U}$ by the polynomial formula.
\begin{enumerate}[label=(\roman*)]
\setlength{\itemsep}{0pt}
\item For $\mathbf{u}\in\mathcal{U}$, $\mathbf{u}=\sum\limits_{\lambda\in\operatorname{spec}(T)}\rho_{\lambda}\mathbf{u}$ by the resolution identity in Theorem \ref{thm:algebraic-spectral-resolution-diagonalizable}(i), while $\rho_{\lambda}\mathbf{u}\in\mathcal{U}_{\lambda}$ by the preceding inclusion and part (ii) of that theorem. Thus $\mathcal{U}=\sum\limits_{\lambda\in\operatorname{spec}(T)}\mathcal{U}_{\lambda}$, and the sum is direct by Proposition \ref{prop:eigenspaces-for-distinct-eigenvalues-form-direct-sum}(i).
\item If $\mathcal{U}=\{\mathbf{0}\}$, the empty basis diagonalizes $T|_{\mathcal{U}}$. Suppose $\mathcal{U}\neq\{\mathbf{0}\}$. Concatenating bases of the nonzero summands forms basis of $\mathcal{U}$ by Theorem \ref{thm:direct-sum-iff-dimensions-add}(ii) and part (i). Its members are eigenvectors of $T|_{\mathcal{U}}$, so restriction is diagonalizable by Theorem \ref{thm:conditions-equivalent-to-diagonalizability}(ii),(i).
\end{enumerate}
\hlt{Structural Commentary}: Polynomial spectral projections preserve a $T$-invariant subspace and split it along the ambient eigenspaces.\\
\hlt{Key Insight}: Restricting a diagonalizable operator restricts each spectral projection as well.
\end{proof}
\end{lemma}

\begin{theorem} \label{thm:simultaneous-diagonalization-commuting-diagonalizable-family} \hlt{Simultaneous Diagonalization of Commuting Diagonalizable Operators}\\
Let $\mathcal{V}$ be a nonzero finite-dimensional vector space, and let $\emptyset\neq\mathcal{F}\subseteq\LL(\mathcal{V})$ be finite with each member diagonalizable.\\
Then $(\forall S,T\in\mathcal{F})\,ST=TS$ iff $(\exists\text{ a basis }\beta\text{ of }\mathcal{V})$ such that $(\forall T\in\mathcal{F})\,[T]_{\beta\leftarrow\beta}$ is diagonal.
\begin{proof}
\begin{enumerate}[label=\roman*.]
\setlength{\itemsep}{0pt}
\item[$(\Rightarrow)$] Assume $(\forall S,T\in\mathcal{F})\,ST=TS$, and proceed by strong induction on $n\defeq\dim(\mathcal{V})$. The case $n=1$ is scalar.\\
Suppose $n>1$ and the result holds in every smaller positive dimension. If each $T\in\mathcal{F}$ is scalar, any basis is common.\\
Otherwise choose a nonscalar $S\in\mathcal{F}$. The decomposition $\mathcal{V}=\bigoplus\limits_{\lambda\in\operatorname{spec}(S)}E_{\lambda}(S)$ holds by Theorem \ref{thm:conditions-equivalent-to-diagonalizability}(i),(iv).\\
If it had one summand, then $\operatorname{spec}(S)=\{\lambda\}$, $S=\lambda I_{\mathcal{V}}$ by right-side identities in Theorem \ref{thm:algebraic-spectral-resolution-diagonalizable}(i), contradicting choice of $S$.\\
Hence $(\forall\lambda\in\operatorname{spec}(S))\,1\leq\dim E_{\lambda}(S)<n$. Fix $\lambda\in\operatorname{spec}(S)$.\\
If $T\in\mathcal{F}$ and $\mathbf{x}\in E_{\lambda}(S)$, then $S(T\mathbf{x})=T(S\mathbf{x})=\lambda T\mathbf{x}$ by commutation. Thus $E_{\lambda}(S)\in\operatorname{Inv}(T)$. Set $\mathcal{F}_{\lambda}\defeq\{T|_{E_{\lambda}(S)}:T\in\mathcal{F}\}$.\\
Every member of family is diagonalizable by Lemma \ref{lem:diagonalizable-restriction-to-invariant-subspace}(ii), and family is finite, nonempty, commuting by construction.\\
A common eigenbasis of $E_{\lambda}(S)$ exists by induction. Concatenating these bases over $\lambda$ forms a basis of $\mathcal{V}$ by Theorem \ref{thm:direct-sum-iff-dimensions-add}(ii), and it diagonalizes each member of $\mathcal{F}$.
\item[$(\Leftarrow)$] Choose a basis $\beta$ satisfying the right side. For $S,T\in\mathcal{F}$, set $D_S=[S]_{\beta\leftarrow\beta}$ and $D_T=[T]_{\beta\leftarrow\beta}$.\\
Since $D_S,D_T$ are diagonal, $[ST]_{\beta\leftarrow\beta}=D_SD_T=D_TD_S=[TS]_{\beta\leftarrow\beta}$ by Proposition \ref{prop:composition-and-matrix-multiplication}. Hence $ST=TS$ by Theorem \ref{thm:linear-maps-and-matrices-same-data}(i).
\end{enumerate}
\hlt{Structural Commentary}: A nonscalar member supplies proper common invariant eigenspaces for induction.\\
\hlt{Key Insight}: Commuting operators preserve the eigenspace decomposition of a diagonalizable operator.
\end{proof}
\end{theorem}

\begin{definition} \label{def:matrix-commutant} \hlt{Commutant}\\
Let $n\geq1$. For $\emptyset\neq\mathcal{F}\subseteq\operatorname{Mat}(n,n,\mathbb{F})$, the \hlt{commutant} of $\mathcal{F}$ is $\operatorname{Comm}(\mathcal{F})\defeq\{X\in\operatorname{Mat}(n,n,\mathbb{F}):(\forall\,A\in\mathcal{F})\,XA=AX\}$.\\
For a single matrix $A$, write $\operatorname{Comm}(A)$.
\end{definition}

\begin{lemma} \label{lem:block-diagonal-commutant-distinct-scalars} \hlt{Commuting with a Block Scalar Diagonal Matrix}\\
Let $k,n_1,\ldots,n_k\geq1$, set $n\defeq\sum\limits_{i=1}^{k}n_i$, let $\lambda_1,\ldots,\lambda_k\in\mathbb{F}$, and set $\Lambda\defeq\lambda_1 I_{n_1}\oplus\cdots\oplus\lambda_k I_{n_k}$.\\
Partition $A=[A_{ij}]\in\operatorname{Mat}(n,n,\mathbb{F})$ conformally with $\Lambda$, so $A_{ij}\in\operatorname{Mat}(n_i,n_j,\mathbb{F})$.
\begin{enumerate}[label=(\roman*)]
\setlength{\itemsep}{0pt}
\item If $A$ is block diagonal, then $A\Lambda=\Lambda A$.
\item If $A\Lambda=\Lambda A$ and $\lambda_1,\ldots,\lambda_k$ are distinct, then $A$ is block diagonal.
\end{enumerate}
\begin{proof}
By block multiplication, $(A\Lambda)_{ij}=\lambda_jA_{ij}$ and $(\Lambda A)_{ij}=\lambda_iA_{ij}$.
\begin{enumerate}[label=(\roman*)]
\setlength{\itemsep}{0pt}
\item If $A$ is block diagonal, then $A_{ij}=0$ for $i\neq j$, while the two diagonal-block expressions coincide. Hence $A\Lambda=\Lambda A$.
\item If $A\Lambda=\Lambda A$, then $(\forall i,j\in\{1,\ldots,k\})\,(\lambda_j-\lambda_i)A_{ij}=0$. Distinctness of $\lambda_i$ implies $A_{ij}=0$ whenever $i\neq j$, so $A$ is block diagonal.
\end{enumerate}
\hlt{Structural Commentary}: Block comparison reduces commutation to scalar off-diagonal equations.\\
\hlt{Key Insight}: Distinct block scalars force every off-diagonal block to vanish.
\end{proof}
\end{lemma}

\begin{example} \hlt{Distinctness Is Necessary}\\
The distinctness hypothesis in Lemma \ref{lem:block-diagonal-commutant-distinct-scalars} cannot be dropped. If $\lambda_1=\cdots=\lambda_k=1$, then $\Lambda=I_n$ and $\operatorname{Comm}(\Lambda)=\operatorname{Mat}(n,n,\mathbb{F})$, which contains matrices that are not block diagonal for the chosen partition.
\end{example}

\begin{theorem} \label{thm:commutant-of-diagonalizable-matrix} \hlt{Commutant of a Diagonalizable Matrix}\\
Let $n\geq1$, $A\in\operatorname{Mat}(n,n,\mathbb{C})$ be diagonalizable. Enumerate distinct eigenvalues as $\mu_1,\ldots,\mu_d$, where $d\geq1$, set $n_i\defeq\dim_{\mathbb{C}}E_{\mu_i}(A)$.\\
Choose $S\in\operatorname{Mat}(n,n,\mathbb{C})^{\times}$ satisfying $A=S(\mu_1I_{n_1}\oplus\cdots\oplus\mu_dI_{n_d})S^{-1}$. The following hold:
\begin{enumerate}[label=(\roman*)]
\setlength{\itemsep}{0pt}
\item For $B\in\operatorname{Mat}(n,n,\mathbb{C})$, $B\in\operatorname{Comm}(A)$ iff $(\exists\,(X_i)_{i=1}^{d}\in\prod\limits_{i=1}^{d}\operatorname{Mat}(n_i,n_i,\mathbb{C}))$ such that $B=S(X_1\oplus\cdots\oplus X_d)S^{-1}$.
\item $\dim_{\mathbb{C}}\operatorname{Comm}(A)=\sum\limits_{i=1}^{d}n_i^2$.
\end{enumerate}
\begin{proof}
Set $\Lambda\defeq\mu_1I_{n_1}\oplus\cdots\oplus\mu_dI_{n_d}$.
\begin{enumerate}[label=(\roman*)]
\setlength{\itemsep}{0pt}
\item[$(i)$] ($\Rightarrow$) Suppose $B\in\operatorname{Comm}(A)$. Then $(S^{-1}BS)\Lambda=\Lambda(S^{-1}BS)$ by multiplying $BA=AB$ on left by $S^{-1}$ and on right by $S$.\\
Hence $S^{-1}BS=X_1\oplus\cdots\oplus X_d$ for suitable $X_i\in\operatorname{Mat}(n_i,n_i,\mathbb{C})$ by Lemma \ref{lem:block-diagonal-commutant-distinct-scalars}(ii).\\
Thus $B=S(X_1\oplus\cdots\oplus X_d)S^{-1}$ by multiplication by $S$ and $S^{-1}$.
\item[$(i)$] ($\Leftarrow$) Suppose $B=S(X_1\oplus\cdots\oplus X_d)S^{-1}$. The block-diagonal matrix $S^{-1}BS$ commutes with $\Lambda$ by Lemma \ref{lem:block-diagonal-commutant-distinct-scalars}(i).\\
Hence $BA=AB$ by conjugating this equality by $S$, and $B\in\operatorname{Comm}(A)$ by Definition \ref{def:matrix-commutant}.
\item[$(ii)$] Conjugation by $S$ is a linear bijection with inverse conjugation by $S^{-1}$.\\
The block-diagonal subspace has a basis formed by embedding the canonical matrix-unit bases of its $d$ diagonal blocks.\\
Its dimension, and therefore that of $\operatorname{Comm}(A)$ by part (i), is $\sum\limits_{i=1}^{d}n_i^2$ by Proposition \ref{prop:matrix-spaces-are-vector-spaces}(ii).
\end{enumerate}
\hlt{Structural Commentary}: Similarity yields block-diagonal form and transfers its dimension.\\
\hlt{Key Insight}: Commuting matrices preserve eigenspaces but act freely within each one.
\end{proof}
\end{theorem}

\begin{theorem} \label{thm:commuting-family-common-eigenvector} \hlt{Common Eigenvector for a Commuting Family}\\
Let $\mathcal{V}$ be a nonzero finite-dimensional complex vector space, and let $\emptyset\neq\mathcal{F}\subseteq\LL(\mathcal{V})$ satisfy $(\forall S,T\in\mathcal{F})\,ST=TS$.\\
Then $(\exists\,\mathbf{v}\in\mathcal{V}\setminus\{\mathbf{0}\})(\exists\,(\lambda_T)_{T\in\mathcal{F}}\in\mathbb{C}^{\mathcal{F}})$ such that $(\forall T\in\mathcal{F})\,T\mathbf{v}=\lambda_T\mathbf{v}$.
\begin{proof}
\begin{enumerate}[label=(\roman*)]
\setlength{\itemsep}{0pt}
\item \hlt{Finite family}: Suppose $\mathcal{F}=\{T_1,\ldots,T_m\}$, and induct on $m\geq1$. The case $m=1$ holds by Theorem \ref{thm:complex-operator-has-eigenvalue}. Suppose $m>1$.\\
Choose $\lambda\in\operatorname{spec}(T_1)$ by that theorem, and set $\mathcal{U}\defeq E_{\lambda}(T_1)\neq\{\mathbf{0}\}$.\\
If $\mathbf{u}\in\mathcal{U}$ and $j\in\{2,\ldots,m\}$, then $(T_1-\lambda I_{\mathcal{V}})T_j\mathbf{u}=T_j(T_1-\lambda I_{\mathcal{V}})\mathbf{u}=\mathbf{0}$ by commutation. Thus $\mathcal{U}$ is invariant under $T_2,\ldots,T_m$.\\
The restrictions to $\mathcal{U}$ commute, so a common eigenvector $\mathbf{0}\neq\mathbf{v}\in\mathcal{U}$ exists by induction.\\
Then $\mathbf{v}\in E_{\lambda}(T_1)$, so it is also an eigenvector of $T_1$.
\item \hlt{Finite reduction}: For an arbitrary family, set $\mathcal{W}\defeq\operatorname{span}(\mathcal{F})$, which is finite-dimensional by Theorem \ref{thm:linear-maps-and-matrices-same-data}(ii).\\
If $\mathcal{W}=\{0\}$, then $(\forall T\in\mathcal{F})\,T=0$ by $\mathcal{F}\subseteq\mathcal{W}$, so any $\mathbf{v}\in\mathcal{V}\setminus\{\mathbf{0}\}$ is common with eigenvalue $0$. Suppose $\mathcal{W}\neq\{0\}$.\\
Choose a basis $(R_1,\ldots,R_q)$ of $\mathcal{W}$ by Corollary \ref{cor:finite-dimensional-spaces-and-subspaces-have-bases}. Each $R_j$ is a finite linear combination of members of $\mathcal{F}$.\\
The operators occurring in these combinations form a finite nonempty set $\mathcal{F}_0\subseteq\mathcal{F}$ with $\operatorname{span}(\mathcal{F}_0)=\mathcal{W}$.
\item \hlt{Spanning extension}: Suppose $\mathcal{W}\neq\{0\}$ and let $\mathcal{F}_0$ be as in part (ii). Choose $\mathbf{v}\neq\mathbf{0}$ and $(\mu_T)_{T\in\mathcal{F}_0}$ satisfying $(\forall T\in\mathcal{F}_0)\,T\mathbf{v}=\mu_T\mathbf{v}$ by part (i). If $S\in\mathcal{F}$, write $S=\sum\limits_{T\in\mathcal{F}_0}a_TT$. Then $S\mathbf{v}=(\sum\limits_{T\in\mathcal{F}_0}a_T\mu_T)\mathbf{v}$.
\end{enumerate}
\hlt{Structural Commentary}: Finite induction handles finite families; finite-dimensionality reduces arbitrary families.\\
\hlt{Key Insight}: Pairwise commutation keeps the remaining operators inside the first operator's eigenspace.
\end{proof}
\end{theorem}

\begin{theorem} \label{thm:simultaneous-schur-triangularization-commuting-family} \hlt{Simultaneous Schur Triangularization}\\
Let $n\geq1$ and let $\emptyset\neq\mathcal{F}\subseteq\operatorname{Mat}(n,n,\mathbb{C})$ satisfy $(\forall A,B\in\mathcal{F})\,AB=BA$. The following hold:
\begin{enumerate}[label=(\roman*)]
\setlength{\itemsep}{0pt}
\item $(\exists\text{ a unitary }U\in\operatorname{Mat}(n,n,\mathbb{C}))$ such that $(\forall A\in\mathcal{F})\,U^*AU$ is upper triangular.
\item If $\mathcal{F}\subseteq\operatorname{Mat}(n,n,\mathbb{R})$, $(\forall A\in\mathcal{F})\,\operatorname{spec}(A)\subseteq\mathbb{R}$, then $(\exists\text{ real orthogonal }U\in\operatorname{Mat}(n,n,\mathbb{R}))$ s.t. $(\forall A\in\mathcal{F})\,U^tAU$ upper triangular.
\end{enumerate}
\begin{proof}
Proceed by strong induction on $n$. Both conclusions hold with $U=I_1$ when $n=1$.\\
Suppose $n>1$ and both conclusions hold in each smaller positive dimension.
\begin{enumerate}[label=(\roman*)]
\setlength{\itemsep}{0pt}
\item A common eigenvector for the operators on $\mathbb{C}^n$ exists by Theorem \ref{thm:commuting-family-common-eigenvector}.\\
Normalize it to $\mathbf{x}$ and extend its column to a unitary matrix $V=[\mathbf{x}\ V_2]$ by Theorem \ref{thm:extend-isometric-columns-to-unitary-matrix}(ii).\\
For $A\in\mathcal{F}$, write $M_A\defeq V^*AV=\left[\begin{smallmatrix}\lambda_A&\mathbf{r}_A\\0&\widehat A\end{smallmatrix}\right]$; the lower-left block is zero by $A\mathbf{x}=\lambda_A\mathbf{x}$.\\
For $A,B\in\mathcal{F}$, $M_AM_B=V^*ABV=V^*BAV=M_BM_A$, so $\widehat A\widehat B=\widehat B\widehat A$ by comparison of the lower-right blocks.\\
Thus $\widehat{\mathcal{F}}\defeq\{\widehat A:A\in\mathcal{F}\}$ is nonempty and commuting. A unitary $W$ simultaneously upper triangularizing its members exists by induction. Set $U\defeq V(I_1\oplus W)$. The matrix $U$ is unitary by Proposition \ref{prop:algebra-of-unitary-matrices}(i)--(ii), and $U^*AU$ is upper triangular for $A\in\mathcal{F}$.
\item Regard $\mathcal{F}$ as commuting family on $\mathbb{C}^n$. $\exists$ $\mathbf{0}\neq\mathbf{z}=\mathbf{u}+i\mathbf{v}$, scalars $\lambda_A$ satisfying $(\forall A\in\mathcal{F})\,A\mathbf{z}=\lambda_A\mathbf{z}$ by Theorem \ref{thm:commuting-family-common-eigenvector}.\\
Each $\lambda_A\in\operatorname{spec}(A)\subseteq\mathbb{R}$ by this eigenvector relation and the spectral hypothesis.\\
Since $A$ is real, $A\mathbf{u}=\lambda_A\mathbf{u}$ and $A\mathbf{v}=\lambda_A\mathbf{v}$ by comparing real and imaginary parts.\\
At least one of $\mathbf{u},\mathbf{v}$ is nonzero; normalize a nonzero member to obtain a common real unit eigenvector $\mathbf{x}$.\\
Extend $\mathbf{x}$ to a real orthogonal matrix $V=[\mathbf{x}\ V_2]$ by Theorem \ref{thm:extend-isometric-columns-to-unitary-matrix}(ii), and write $M_A\defeq V^tAV=\left[\begin{smallmatrix}\lambda_A&\mathbf{r}_A\\0&\widehat A\end{smallmatrix}\right]$.\\
The lower-right blocks are real and commute by the calculation in part (i).\\
Fix $A\in\mathcal{F}$ and $\mu\in\operatorname{spec}(\widehat A)$, and choose $\mathbf{0}\neq\mathbf{y}\in\mathbb{C}^{n-1}$ with $\widehat A\mathbf{y}=\mu\mathbf{y}$. If $\mu=\lambda_A$, then $\mu\in\mathbb{R}$. If $\mu\neq\lambda_A$, set $c\defeq\frac{\mathbf{r}_A\mathbf{y}}{\mu-\lambda_A}$.\\
Then $M_A(c,\mathbf{y})=(\mu c,\mu\mathbf{y})=\mu(c,\mathbf{y})$, so $V(c,\mathbf{y})$ is an eigenvector of $A$ for $\mu$. Hence $\mu\in\mathbb{R}$ by hypothesis.\\
Thus $\operatorname{spec}(\widehat A)\subseteq\mathbb{R}$ for $A\in\mathcal{F}$. A real orthogonal $W$ simultaneously upper triangularizing $\widehat{\mathcal{F}}$ exists by induction.\\
Then $U\defeq V(I_1\oplus W)$ is real orthogonal by Proposition \ref{prop:algebra-of-unitary-matrices}(i)--(ii), and $U^tAU$ is upper triangular for $A\in\mathcal{F}$.
\end{enumerate}
\hlt{Structural Commentary}: A common eigenline creates the first block; commutation descends to the compression, enabling induction.\\
\hlt{Key Insight}: Real spectra turn a complex common eigenvector into a real one.
\end{proof}
\end{theorem}

Figure \ref{fig:simultaneous-triangularization-commuting} summarizes the complex induction in Theorem \ref{thm:simultaneous-schur-triangularization-commuting-family}: a common eigenline creates the first block, and pairwise commutation passes to the lower-right blocks. The theorem, rather than the diagram, owns the simultaneous triangularization conclusion.

\begin{illustrationfigure}[H]
  \centering
  \begin{tikzpicture}[
  >=Stealth,
  every node/.style={font=\small},
  box/.style={draw=gray!55, fill=gray!5, rounded corners=2pt, align=center, inner xsep=5pt, inner ysep=4pt},
  hi/.style={draw=blue!60!black, fill=blue!7, rounded corners=2pt, align=center, inner xsep=5pt, inner ysep=4pt},
  arr/.style={->, line width=0.75pt, blue!70!black}
]
  \node[hi] (common) at (0,0)
    {$\mathbb{C}\mathbf{x}$};
  \node[box] (block) at (3.18,0)
    {$V^*AV=\begin{bmatrix}\lambda_A&*\\0&A'\end{bmatrix}$};
  \node[box] (comm) at (6.60,0)
    {$A'B'=B'A'$};

  \draw[arr] (common.east) -- (block.west);
  \draw[arr] (block.east) -- (comm.west);
\end{tikzpicture}

  \caption{A shared eigenline determines a simultaneous first block; lower-right compressions remain commuting.}
  \label{fig:simultaneous-triangularization-commuting}
\end{illustrationfigure}

\begin{corollary} \label{cor:commuting-sum-product-eigenvalues} \hlt{Eigenvalues of Sums and Products of Commuting Matrices}\\
Let $n\geq1$ and let $A,B\in\operatorname{Mat}(n,n,\mathbb{C})$ satisfy $AB=BA$. Then $(\exists\,\lambda_1,\ldots,\lambda_n,\mu_1,\ldots,\mu_n\in\mathbb{C})$ such that the following hold:
\begin{enumerate}[label=(\roman*)]
\setlength{\itemsep}{0pt}
\item $\operatorname{spec}(A)=\{\lambda_1,\ldots,\lambda_n\}$ and $\operatorname{spec}(B)=\{\mu_1,\ldots,\mu_n\}$, with all $n$ diagonal occurrences retained in each list;
\item $\operatorname{spec}(A+B)=\{\lambda_i+\mu_i:i\in\{1,\ldots,n\}\}$;
\item $\operatorname{spec}(AB)=\{\lambda_i\mu_i:i\in\{1,\ldots,n\}\}$.
\end{enumerate}
\begin{proof}
Choose a unitary $U$ by Theorem \ref{thm:simultaneous-schur-triangularization-commuting-family}(i), and set $T\defeq U^*AU$, $R\defeq U^*BU$, $\lambda_i\defeq t_{ii}$, and $\mu_i\defeq r_{ii}$.\\
Fix $M\in\operatorname{Mat}(n,n,\mathbb{C})$ and $\theta\in\mathbb{C}$. For $\mathbf{x}\neq\mathbf{0}$, $(U^*MU)\mathbf{x}=\theta\mathbf{x}$ iff $M(U\mathbf{x})=\theta(U\mathbf{x})$, and $U\mathbf{x}\neq\mathbf{0}$. Hence $\operatorname{spec}(U^*MU)=\operatorname{spec}(M)$.
\begin{enumerate}[label=(\roman*)]
\setlength{\itemsep}{0pt}
\item The matrices $T$ and $R$ are upper triangular with diagonals $(\lambda_i)_{i=1}^{n}$ and $(\mu_i)_{i=1}^{n}$. Their spectra are the underlying sets of these lists by Proposition \ref{prop:eigenvalues-of-upper-triangular-operator}. Hence the identities in part (i) hold by the preceding similarity calculation.
\item The matrix $T+R=U^*(A+B)U$ is upper triangular with diagonal $(\lambda_i+\mu_i)_{i=1}^{n}$.\\
Its spectrum is the asserted set by Proposition \ref{prop:eigenvalues-of-upper-triangular-operator}. Hence the identity in part (ii) holds by the preceding similarity calculation.
\item The matrix $TR=U^*ABU$ is upper triangular with diagonal $(\lambda_i\mu_i)_{i=1}^{n}$.\\
Its spectrum is the asserted set by Proposition \ref{prop:eigenvalues-of-upper-triangular-operator}. Hence identity in part (iii) holds by the preceding similarity calculation.
\end{enumerate}
\hlt{Structural Commentary}: Common triangular coordinates reduce both spectra to diagonal arithmetic.\\
\hlt{Key Insight}: A common diagonal pairing controls both resulting spectra.
\end{proof}
\end{corollary}
